\documentclass[10pt,a4paper]{amsart}
\usepackage[margin=19mm]{geometry}
\usepackage{amsmath,amssymb,mathtools}
\usepackage[scr=rsfs,cal=euler]{mathalfa}
\usepackage{xfrac}
\usepackage[unicode,hidelinks,bookmarks=false]{hyperref}
\newcommand{\ie}{i.e.\ }
\newcommand{\cf}{cf.\ }
\newcommand{\resp}{respectively\ }
\newcommand{\Subsection}{\subsection}
\providecommand{\nobreakdash}{\nobreak\hbox{-}}
\setlength{\emergencystretch}{2em}
\multlinegap0pt
\usepackage{amssymb}
\usepackage{xcolor}
\usepackage[normalem]{ulem}
\newtheorem{lemma}{Lemma}
\title{Positive Logarithmic Hausdorff Measures of Exceptional Sets for the $p$-adic and $t$-adic Littlewood Conjectures}
\author{Dzmitry Badziahin, Volodymyr Pavlenkov, Evgeniy Zorin}
\date{Source: arXiv:2608.22078v1 (CC BY 4.0)}
\begin{document}
\maketitle

\noindent\emph{Source and licence.} Positive Logarithmic Hausdorff Measures of Exceptional Sets for the $p$-adic and $t$-adic Littlewood Conjectures. Version arXiv:2608.22078v1 (CC BY 4.0), distributed by arXiv under the Creative Commons Attribution 4.0 International licence, http://creativecommons.org/licenses/by/4.0/, as recorded in the arXiv licence metadata for this exact version and shown on the abstract page https://arxiv.org/abs/2608.22078v1. Full licence notice: \texttt{licenses/arxiv-2608.22078-CC-BY-4.0.txt}.\par\smallskip

\noindent\textit{Editorial scope.} Counted unit: the article's lemma un\_contrex\_lemma\_tadic (source0.tex lines 868--875). The article's complete original proof of it is delivered with it (source0.tex lines 877--1172, one \texttt{proof} environment of 296 lines). No mathematics is written, repaired or re-derived: the statement and the proof are the article's own lines, in the article's own notation, and no retained line is substituted or re-flowed.  Retained uncounted as the source-local context the proof needs: lines 792--796; lines 807--810; lines 862--866.  Nothing was omitted from any retained span, so the package carries no omission record.  No retained line contains a citation, so the wrapper carries no bibliography and no bibliography entry.  No retained line contains a figure environment, an image-inclusion command or drawing code, so no image is delivered and the package declares no asset dependency.  Editorial formatting only, in addition: an AMS \texttt{amsart} preamble that restates the article's own declarations and macros used by the retained lines, namely the environments \texttt{lemma} on separate counters, together with the standard packages those lines need.  The retained environments print in this document as Lemma 1 in order of appearance, whereas the article prints the same items with the numbering of its own full document.  The complete native source of the article as distributed in the arXiv e-print for this exact version is bundled unchanged as \texttt{sources/208301/source0.tex}.
\begin{equation} \label{eq:d-and-Z}
 d=\frac{q-1}{2},
 \qquad
 Z=\{\zeta\in\mathbb F_q^\times:\zeta^d=-1\}.
\end{equation}

\begin{equation} \label{eq:parameter-space}
 \mathbf a=(a_{k,\zeta})_{k\geq0,\,\zeta\in Z},
 \qquad a_{k,\zeta}\in\mathbb F_q^\times,
\end{equation}

\begin{equation} \label{eq:parametrised-functional-equation}
 \Lambda_{\mathbf a}(t)
 =\Lambda_{\sigma\mathbf a}(t^2)
  +\sum_{\zeta\in Z}a_{0,\zeta}\frac{t}{t^2-\zeta}.
\end{equation}

\begin{lemma}[Uniform Counterexample Lemma]\label{un_contrex_lemma_tadic}
For every $\mathbf a$ satisfying \eqref{eq:parameter-space},
\begin{equation} \label{eq:uniform-parametrised-bound}
 \inf_{\substack{0\ne M\in\mathbb F_q[t]\\m\geq0}}
 |M|\,|\langle Mt^m\Lambda_{\mathbf a}\rangle|
 \geq q^{-2(q-1)}.
\end{equation}
\end{lemma}

\begin{proof}[Proof of Lemma~\ref{un_contrex_lemma_tadic}]
Put
\begin{equation} \label{eq:D-polynomial}
 D(t)=t^{q-1}-1=t^{2d}-1=(t^d-1)(t^d+1).
\end{equation}
We first prove that there are no $\mathbf a$, $m\geq0$ and non-zero
$Q\in\mathbb F_q[t]$ such that
\begin{equation} \label{eq:key-violation}
 D\mid Q
 \qquad\text{and}\qquad
 |Q|\,|\langle Qt^m\Lambda_{\mathbf a}\rangle|<1.
\end{equation}
Suppose otherwise, and choose such a triple $(\mathbf a, m, Q)$ with
$\deg Q$ minimal among all parameter collections $\mathbf a$. Write
\begin{equation} \label{eq:Q-parity-decomposition}
 Q(t)=D(t)\widetilde Q(t),
 \qquad
 \widetilde Q(t)=\widetilde Q_0(t^2)+t\widetilde Q_1(t^2),
\end{equation}
and set
\begin{equation} \label{eq:Qi-definition}
 Q_i(t)=(t^d-1)\widetilde Q_i(t),\qquad i=0,1.
\end{equation}
Then
\begin{equation} \label{eq:Q-even-odd}
 Q(t)=Q_0(t^2)+tQ_1(t^2).
\end{equation}
Write
\begin{equation} \label{eq:degree-parities}
 \deg Q=2s+\varepsilon,
 \qquad
 m=2n+\delta,
 \qquad \varepsilon,\delta\in\{0,1\}.
\end{equation}
Since $D\mid Q$, we have $s\geq d$.  Moreover,
\begin{align} \label{eq:parity-degrees}
 \deg Q_\varepsilon&=s,
 &\deg\widetilde Q_\varepsilon&=s-d,\notag\\
 \deg Q_{1-\varepsilon}&\leq s-1+\varepsilon,
 &\deg\widetilde Q_{1-\varepsilon}&\leq s-d-1+\varepsilon.
\end{align}
In particular, $Q_\varepsilon\ne0$ and $\deg Q_\varepsilon<\deg Q$.

One can easily verify that
$$
 \frac{t^{2n+1+\delta}}{t^2-\zeta} -
 \frac{\zeta^{n+\delta}t^{1-\delta}}{t^2-\zeta} \in \mathbb F_q[t],
$$
therefore
\[
 \left\langle
 \frac{t^{2n+1+\delta}}{t^2-\zeta}
 \right\rangle
 =
 \frac{\zeta^{n+\delta}t^{1-\delta}}{t^2-\zeta}.
\]
We apply this formula to get
$$
\left\langle \frac{(Q_0(t^2) + tQ_1(t^2))t}{t^2-\zeta}\right\rangle
= \frac{tQ_0(\zeta) + \zeta Q_1(\zeta)}{t^2-\zeta}.
$$
By using this equation and separating the even and odd powers of
$t$, we obtain
\begin{align*}
 \left\langle Qt^m\Lambda_{\mathbf a}\right\rangle
 &=
 \left\langle
 Q_\delta(t^2)t^{2n+2\delta}
 \Lambda_{\sigma\mathbf a}(t^2)
 +
 \sum_{\zeta\in Z}
 a_{0,\zeta}
 \frac{\zeta^{n+1}Q_{1-\delta}(\zeta)}
      {t^2-\zeta}
 \right\rangle
 \\
 &\quad
 +t\left\langle
 Q_{1-\delta}(t^2)t^{2n}
 \Lambda_{\sigma\mathbf a}(t^2)
 +
 \sum_{\zeta\in Z}
 a_{0,\zeta}
 \frac{\zeta^{n+\delta}Q_\delta(\zeta)}
      {t^2-\zeta}
 \right\rangle .
\end{align*}



%Therefore, multiplying \eqref{eq:parametrised-functional-equation}
%by $t^{2n+\delta}$ gives
%\begin{equation} \label{eq:shifted-functional-equation}
% \left\langle t^{2n+\delta}\Lambda_{\mathbf a}(t)\right\rangle
% =
% \left\langle
% t^{2n+\delta}\Lambda_{\sigma\mathbf a}(t^2)
% \right\rangle
% +
% \sum_{\zeta\in Z}
% a_{0,\zeta}
% \frac{\zeta^{n+\delta}t^{1-\delta}}{t^2-\zeta}.
%\end{equation}

For a Laurent series $R$, we write $R=O(t^u)$ if $\deg R\leq u$. The
strict inequality in \eqref{eq:key-violation}, together with the
discreteness of the absolute value, implies that $\langle
Qt^m\Lambda_{\mathbf a}\rangle =O(t^{-\deg Q-1})$. Consequently,
\begin{align}
 \left\langle
 Q_\delta(t)t^{n+\delta}\Lambda_{\sigma\mathbf a}(t)
 +
 \sum_{\zeta\in Z}
 a_{0,\zeta}
 \frac{\zeta^{n+1}Q_{1-\delta}(\zeta)}
      {t-\zeta}
 \right\rangle
 &=O(t^{-s-1}),
 \label{eq:descent-first}
 \\
 \left\langle
 Q_{1-\delta}(t)t^n\Lambda_{\sigma\mathbf a}(t)
 +
 \sum_{\zeta\in Z}
 a_{0,\zeta}
 \frac{\zeta^{n+\delta}Q_\delta(\zeta)}
      {t-\zeta}
 \right\rangle
 &=O(t^{-s-1-\varepsilon}).
 \label{eq:descent-second}
\end{align}
These are the two basic descent identities.

For $\zeta\in Z$, equations \eqref{eq:d-and-Z} and
\eqref{eq:Qi-definition} give
\begin{equation} \label{eq:Qi-at-zeta}
 Q_i(\zeta)=-2\widetilde Q_i(\zeta).
\end{equation}
The degree bounds in \eqref{eq:parity-degrees} imply
\[
 \deg\widetilde Q_{1-\delta}\leq s-d,
 \qquad
 \deg\bigl(t^\delta\widetilde Q_\delta\bigr)
 \leq s-d+\varepsilon\delta.
\]
Hence, one can easily check that after multiplying
\eqref{eq:descent-first} by $\widetilde Q_{1-\delta}(t)$
%resulting error is
%\[
% O(t^{-s-1+s-d})=O(t^{-d-1}),
%\]
and multiplying \eqref{eq:descent-second} by $t^\delta\widetilde
Q_\delta(t)$ both expressions are of order $O(t^{-d-1})$.
%\[
% O(t^{-s-1-\varepsilon+s-d+\varepsilon\delta})
% =
% O(t^{-d-1-\varepsilon(1-\delta)})
% =
% O(t^{-d-1}).
%\]

After each multiplication we take fractional parts.  We use the identity
\[
 \left\langle\frac{P(t)}{t-\zeta}\right\rangle
 =
 \frac{P(\zeta)}{t-\zeta},
 \qquad P\in\mathbb F_q[t].
\]
Together with \eqref{eq:Qi-at-zeta}, this gives
$$
\left\langle (t^d-1)\widetilde Q_0(t)\widetilde Q_1(t)
 t^{n+\delta}\Lambda_{\sigma\mathbf a}(t) - 2 \sum_{\zeta\in Z}
 a_{0,\zeta}
 \frac{\zeta^{n+1}\widetilde Q_{1-\delta}(\zeta)^2}
      {t-\zeta}\right\rangle = O(t^{-d-1}),
$$
$$
\left\langle (t^d-1)\widetilde Q_0(t)\widetilde Q_1(t)
 t^{n+\delta}\Lambda_{\sigma\mathbf a}(t) - 2\sum_{\zeta\in Z}
 a_{0,\zeta}
 \frac{\zeta^{n+2\delta}\widetilde Q_\delta(\zeta)^2}
      {t-\zeta}
 \right\rangle  =O(t^{-d-1}).
$$
%
%\begin{align*}
% \Biggl\langle
% &(t^d-1)\widetilde Q_0(t)\widetilde Q_1(t)
% t^{n+\delta}\Lambda_{\sigma\mathbf a}(t)
% \\
% &\quad
% -2\sum_{\zeta\in Z}
% a_{0,\zeta}
% \frac{\zeta^{n+1}\widetilde Q_{1-\delta}(\zeta)^2}
%      {t-\zeta}
% \Biggr\rangle
% =O(t^{-d-1})
%\end{align*}
%and
%\begin{align*}
% \Biggl\langle
% &(t^d-1)\widetilde Q_0(t)\widetilde Q_1(t)
% t^{n+\delta}\Lambda_{\sigma\mathbf a}(t)
% \\
% &\quad
% -2\sum_{\zeta\in Z}
% a_{0,\zeta}
% \frac{\zeta^{n+2\delta}\widetilde Q_\delta(\zeta)^2}
%      {t-\zeta}
% \Biggr\rangle
% =O(t^{-d-1}).
%\end{align*}
By subtracting these two relations and dividing by the non-zero
scalar $2(-1)^\delta$, we obtain
\begin{equation} \label{eq:Vandermonde-input}
 \sum_{\zeta\in Z}
 a_{0,\zeta}\zeta^{n+\delta}
 \frac{\widetilde Q_0(\zeta)^2
       -\zeta\widetilde Q_1(\zeta)^2}
      {t-\zeta}
 =O(t^{-d-1}).
\end{equation}
Since
\[
 \frac1{t-\zeta}=\sum_{j\geq1}\zeta^{j-1}t^{-j},
\]
the first $d$ coefficients in \eqref{eq:Vandermonde-input} give
\begin{equation} \label{eq:Vandermonde-system}
 \sum_{\zeta\in Z}
 a_{0,\zeta}\zeta^{n+\delta+j-1}
 \bigl(\widetilde Q_0(\zeta)^2
       -\zeta\widetilde Q_1(\zeta)^2\bigr)=0,
 \qquad 1\leq j\leq d.
\end{equation}
We look at them as a system of linear equations in variables
$\widetilde{Q}_0(\zeta)^2 - \zeta\widetilde{Q}_1(\zeta)^2$ with the
corresponding matrix $(a_{0,\zeta}\zeta^{n+\delta+j-1})_{\zeta\in Z,
1\le j\le d}$. Its determinant equals
$$
\prod_{\zeta\in Z} a_{0,\zeta}\zeta^{n+\delta} \cdot \det
(\zeta^{j-1})_{\zeta\in Z, 1\le j\le d }\neq 0,
$$
because the matrix on the right hand side is the Vandermonde matrix
which is invertible. Therefore it has only a zero solution or in
other words
%The Vandermonde matrix $(\zeta^{j-1})_{\zeta\in Z,1\leq
%j\leq d}$ is invertible, and all $a_{0,\zeta}\zeta^{n+\delta}$ are
%non-zero. Hence
$$
 \widetilde Q_0(\zeta)^2
 =\zeta\widetilde Q_1(\zeta)^2
 \qquad(\zeta\in Z).
$$
Every $\zeta\in Z$ is a quadratic non-residue. The last equation
therefore forces
\begin{equation} \label{eq:both-vanish}
 \widetilde Q_0(\zeta)=\widetilde Q_1(\zeta)=0
 \qquad(\zeta\in Z).
\end{equation}

The elements of $Z$ are precisely the $d$ roots of $t^d+1$ in
$\mathbb F_q$.  These roots are distinct, since the characteristic of
$\mathbb F_q$ does not divide $d=(q-1)/2$.  Hence
\[
 t^d+1=\prod_{\zeta\in Z}(t-\zeta).
\]
We now deduce from \eqref{eq:both-vanish} that $t^d+1$ divides both
$\widetilde Q_0$ and $\widetilde Q_1$.  Together with
\eqref{eq:Qi-definition}, this gives
\begin{equation} \label{eq:D-divides-Qi}
 D\mid Q_0
 \qquad\text{and}\qquad
 D\mid Q_1.
\end{equation}

If $\varepsilon=\delta$, use \eqref{eq:descent-first}; otherwise use
\eqref{eq:descent-second}.  In either case \eqref{eq:both-vanish}
eliminates the rational sum and gives
\begin{equation} \label{eq:smaller-violation}
 \left|Q_\varepsilon\right|
 \left|\left\langle
 Q_\varepsilon t^{n+\varepsilon\delta}
 \Lambda_{\sigma\mathbf a}\right\rangle\right|<1.
\end{equation}
By \eqref{eq:D-divides-Qi}, the polynomial $Q_\varepsilon$ is
divisible by $D$, whereas \eqref{eq:parity-degrees} gives $0\leq\deg
Q_\varepsilon<\deg Q$.  This contradicts the global minimality of
$\deg Q$.  Thus \eqref{eq:key-violation} is impossible.

Finally, if \eqref{eq:uniform-parametrised-bound} failed, there
would be $M\ne0$ and $m\geq0$ for which the product in that formula
was smaller than $q^{-2(q-1)}=q^{-4d}$.  Taking $Q=DM$,
multiplication by $D$, whose degree is $2d$, can increase each of
the two absolute values by at most $q^{2d}$. We would obtain
\eqref{eq:key-violation}, a contradiction.  This proves the lemma.
\end{proof}

\end{document}
